\section{Derivations} \label{sec:appendix_a}

\subsection{Derivation of the Canonical Investment Model Solution}
\label{app:canonical-investment-derivation}

This section derives the optimal investment condition for the firm described in
Section~\ref{sec:investment_problem}. The firm chooses investment $i_{j,t}$ and
next-period debt $b_{j,t+1}$ to maximize discounted dividends subject to the
capital accumulation equation and the collateral borrowing constraint.
Let $q_{j,t}$ denote the multiplier on the capital accumulation equation and
$\mu_{j,t}$ the multiplier on the collateral constraint. The Lagrangian is

\begin{align}
\mathcal{L} = \sum_{t=0}^{\infty} \beta^t \Big[
& f(z_{j,t},k_{j,t}) - i_{j,t}
- \frac{\phi}{2}\left(\frac{i_{j,t}}{k_{j,t}}\right)^2 k_{j,t}
+ b_{j,t+1} - (1+r)b_{j,t} \\
& + q_{j,t}\big((1-\delta)k_{j,t} + i_{j,t} - k_{j,t+1}\big) \\
& + \mu_{j,t}\big(\theta k_{j,t+1} - b_{j,t+1}\big)
\Big].
\end{align}

The first-order condition with respect to investment is

\begin{equation}
-1 - \phi \frac{i_{j,t}}{k_{j,t}} + q_{j,t} = 0 .
\end{equation}

Rearranging gives the Tobin's $q$ investment condition

\begin{equation}
q_{j,t} = 1 + \phi \frac{i_{j,t}}{k_{j,t}} .
\end{equation}

Solving for the investment rate,

\begin{equation}
\frac{i_{j,t}}{k_{j,t}} = \frac{1}{\phi}(q_{j,t}-1).
\end{equation}

Under quadratic adjustment costs, the investment-capital ratio is linear in marginal $q$.

The first-order condition with respect to next-period debt is

\begin{equation}
1 - \beta(1+r) - \mu_{j,t} = 0 .
\end{equation}

The first-order condition with respect to next-period capital is

\begin{equation}
-q_{j,t} + \beta \mathbb{E}_t \Big[
f_k(z_{j,t+1},k_{j,t+1})
+ (1-\delta) q_{j,t+1}
+ \theta \mu_{j,t+1}
\Big] = 0 .
\end{equation}

Rearranging yields the Euler equation for marginal $q$,

\begin{equation}
q_{j,t} =
\beta \mathbb{E}_t \Big[
f_k(z_{j,t+1},k_{j,t+1})
+ (1-\delta) q_{j,t+1}
+ \theta \mu_{j,t+1}
\Big].
\end{equation}

The borrowing constraint satisfies the complementary slackness conditions

\begin{align}
b_{j,t+1} &\le \theta k_{j,t+1} \\
\mu_{j,t} &\ge 0 \\
\mu_{j,t}\big(\theta k_{j,t+1} - b_{j,t+1}\big) &= 0 .
\end{align}

When the constraint is slack ($\mu_{j,t}=0$), the model reduces to the standard
Tobin's $q$ investment condition. When the constraint binds ($\mu_{j,t}>0$),
capital carries an additional value because it relaxes the borrowing limit.

\subsection{Derivation of Optimal Action Policy \ref{prop:optimal_policy}}

The agent maximizes expected payoff subject to information and probability constraints:
\begin{align*}
    &\max_{\pi_t(\cdot|s_t)} \sum_c \hat{Q}_t(c, s_t) \pi_t(c|s_t) && (\text{Objective Function}) \\
    \text{s.t. } &-\sum_c \pi_t(c|s_t) \ln \pi_t(c|s_t) \geq \kappa \sum_c \sigma_t(c|s_t) && (\text{Information Constraint}) \\
    &\sum_c \pi_t(c|s_t) = 1 && (\text{Probability Constraint})
\end{align*}

The Lagrangian and subsequent First-Order Conditions (FOC) are:
\begin{align*}
    \mathcal{L} &= \sum_c \hat{Q}_t(c, s_t) \pi_t(c|s_t) + \delta_t \left[ -\sum_c \pi_t \ln \pi_t - \kappa \sum \sigma \right] + \mu \left[ \sum_c \pi_t - 1 \right] && (\text{Lagrangian}) \\
    \frac{\partial \mathcal{L}}{\partial \pi_t} &= \hat{Q}_t(c, s_t) - \delta_t \ln \pi_t(c|s_t) - \delta_t + \mu = 0 && (\text{FOC w.r.t. } \pi_t) \\
    &\implies \ln \pi_t(c|s_t) = \frac{\mu - \delta_t}{\delta_t} + \frac{\hat{Q}_t(c, s_t)}{\delta_t} && (\text{Solve for } \ln \pi_t)
\end{align*}

Solving for the optimal policy $\pi_t(c|s_t)$:
\begin{align*}
    \pi_t(c|s_t) &= \exp \left( \frac{\mu - \delta_t}{\delta_t} \right) \exp \left( \frac{\hat{Q}_t(c, s_t)}{\delta_t} \right) && (\text{Exponentiate}) \\
    &= A \cdot \exp \left( \frac{\hat{Q}_t(c, s_t)}{\delta_t} \right) && (\text{Let } A = e^{\frac{\mu-\delta_t}{\delta_t}}) \\
    &= \frac{\exp \left( \frac{\hat{Q}_t(c, s_t)}{\delta_t} \right)}{\sum_{c'} \exp \left( \frac{\hat{Q}_t(c', s_t)}{\delta_t} \right)} && (\text{Normalize via } \sum \pi_t = 1)
\end{align*}

In the limit where the shadow price of attention $\delta_t \to 0$:
\begin{equation*}
    \lim_{\delta_t \to 0} \pi_t(c|s_t) = 
    \begin{cases} 
    1 & \text{if } c = \arg\max_{c'} \hat{Q}_t(c', s_t) \\ 
    0 & \text{otherwise} 
    \end{cases}
\end{equation*}

\subsection{Gaussian Process Regression \& Experience Learning} \label{app:gp_regression_equivalence}

The primal update in \ref{def:experience_learning_update} produces
three objects: $\hat{Q}^E_t$, $\hat{\Sigma}^E_t$, and $\alpha^E_t$. Each
of these objects is a function over the continuous domain $\mathcal{A} \times \mathcal{S}$.
To implement the primal form directly, one would need to evaluate and store
$\hat{Q}^E_t(\mathbf{a}, \mathbf{s})$ at every point in the action-state space,
which is uncountably infinite. In the code, we instead immplement a GP regression which
sidesteps this issue by storing not the updated functions themselves, but the finite collection
of data from which those functions can be reconstructed at any query point on demand.

\paragraph{From $t = 1$.} In the first period, the agent has some prior $Q^*_0 \sim GP(\mu, k)$
they start in state $\mathbf{s}_0$ and land in state $\mathbf{s}_1$ after taking actions $\mathbf{a}_0$.
They observe a noisy signal, their utility from that action in that state, which is 
a noisy observation of the linear functional $Q^*(\mathbf{a}_0, \mathbf{s}_0) - Q^*(\mathbf{a}_1, \mathbf{s}_1) = \eta^E_1 - \varepsilon_1$.
The noise comes from the stochastic transition. As a result, $\delta_1$, the temporal difference
error, is a noisy signal of the linear functional $Q^*(\mathbf{a}_0, \mathbf{s}_0) - \beta Q^*(\mathbf{a}_1, \mathbf{s}_1)$.
By the linearity of covariance.
\begin{align}
    \operatorname{Cov}(\delta_1, Q^*(\mathbf{a}, \mathbf{s})) &= k((\mathbf{a}_0, \mathbf{s}_0), (\mathbf{a}, \mathbf{s})) - \beta k((\mathbf{a}_1, \mathbf{s}_1), (\mathbf{a}, \mathbf{s})) \\
    \operatorname{Var}(\delta_1) &= k((\mathbf{a}_0, \mathbf{s}_0), (\mathbf{a}_0, \mathbf{s}_0)) + \beta^2 k((\mathbf{a}_1, \mathbf{s}_1), (\mathbf{a}_1, \mathbf{s}_1)) - 2\beta k((\mathbf{a}_0, \mathbf{s}_0), (\mathbf{a}_1, \mathbf{s}_1))
\end{align}
Denote $w \equiv \operatorname{Var}(\delta_1)$, the denominator above.
The posterior mean at any query point $(\mathbf{a}, \mathbf{s})$ then follows directly 
from \ref{def:experience_learning_update}:
\begin{equation}
    \hat{Q}^E_1(\mathbf{a}, \mathbf{s}) 
    = \mu(\mathbf{a}, \mathbf{s}) 
    + \alpha^E_1(\mathbf{a}, \mathbf{s}) \cdot \delta_1
    = \mu(\mathbf{a}, \mathbf{s}) 
    + \frac{k((\mathbf{a}_0, \mathbf{s}_0), (\mathbf{a}, \mathbf{s})) 
    - \beta\, k((\mathbf{a}_1, \mathbf{s}_1), (\mathbf{a}, \mathbf{s}))}{w} 
    \cdot \delta_1
\end{equation}
where $\delta_1 = \eta^E_1 + \beta\, \mu(\mathbf{a}_1, \mathbf{s}_1) - \mu(\mathbf{a}_0, \mathbf{s}_0)$
is the temporal difference error evaluated against the prior mean.


After $N$ observations, the dual form maintains four objects. First, two
$N \times d$ matrices of support points: the \emph{decision points}
$\mathbf{X}^{\mathrm{dec}} = \{(\mathbf{a}_0, \mathbf{s}_0), \ldots,
(\mathbf{a}_{N-1}, \mathbf{s}_{N-1})\}$, where actions were taken, and
the \emph{outcome points}
$\mathbf{X}^{\mathrm{out}} = \{(\mathbf{a}_1, \mathbf{s}_1), \ldots,
(\mathbf{a}_N, \mathbf{s}_N)\}$, where the agent landed. Second, the
$N$-vector of observed utilities
$\mathbf{Y} = (\eta^E_1, \ldots, \eta^E_N)^\top$. Third, an $N \times N$
matrix $C^{-1}$, the inverse of the Gram matrix $C$.

\paragraph{Gram matrix.}
The Gram matrix $C$ is the $N \times N$ matrix whose $(i,j)$ entry is the
prior covariance between the $i$-th and $j$-th TD functionals:
\begin{equation}
    C_{ij} = \operatorname{Cov}(\delta_i, \delta_j)
    = k(\mathbf{x}^{\mathrm{dec}}_i, \mathbf{x}^{\mathrm{dec}}_j)
    - \beta\, k(\mathbf{x}^{\mathrm{dec}}_i, \mathbf{x}^{\mathrm{out}}_j)
    - \beta\, k(\mathbf{x}^{\mathrm{out}}_i, \mathbf{x}^{\mathrm{dec}}_j)
    + \beta^2\, k(\mathbf{x}^{\mathrm{out}}_i, \mathbf{x}^{\mathrm{out}}_j)
\end{equation}
The four-term expansion follows from the bilinearity of covariance applied to
the linear functional $\delta_i = Q^*(\mathbf{x}^{\mathrm{dec}}_i) -
\beta\, Q^*(\mathbf{x}^{\mathrm{out}}_i) + \varepsilon_i$. Expanding
$\operatorname{Cov}(Q^*(\mathbf{x}^{\mathrm{dec}}_i) - \beta\,
Q^*(\mathbf{x}^{\mathrm{out}}_i),\; Q^*(\mathbf{x}^{\mathrm{dec}}_j) -
\beta\, Q^*(\mathbf{x}^{\mathrm{out}}_j))$ and using
$\operatorname{Cov}(Q^*(\mathbf{x}), Q^*(\mathbf{x}')) = k(\mathbf{x},
\mathbf{x}')$ yields the four kernel evaluations above. Intuitively,
$C_{ij}$ measures how much the $i$-th and $j$-th temporal difference
observations tell us about one another under the prior.

\paragraph{Prior TD means.}
Define the $N$-vector of prior TD means $\mathbf{M}$, whose $i$-th entry
is the prior expectation of the $i$-th TD functional:
\begin{equation}
    M_i = \mu(\mathbf{a}_{i-1}, \mathbf{s}_{i-1})
    - \beta\, \mu(\mathbf{a}_i, \mathbf{s}_i)
\end{equation}
The residual $\mathbf{Y} - \mathbf{M}$ is then the $N$-vector whose $i$-th entry is
$\eta^E_i - M_i = \eta^E_i - \mu(\mathbf{a}_{i-1}, \mathbf{s}_{i-1})
+ \beta\, \mu(\mathbf{a}_i, \mathbf{s}_i) = \delta_i$, the temporal
difference error of the $i$-th observation evaluated against the
\emph{original prior} $\mu$. In the primal (Kalman) form of \ref{def:experience_learning_update}, the
TD error at step $t$ is evaluated against the current posterior
mean $\hat{Q}^E_{t-1}$, which has already absorbed all previous
observations. In the dual form, every residual is evaluated against
the same unchanging prior $\mu$. The two forms nonetheless yield
identical posteriors because the inverse Gram matrix $C^{-1}$
compensates for this difference.

\paragraph{Decorrelation.}
To reconstruct the posterior at any query point $(\mathbf{a}, \mathbf{s})$,
we compute the $N$-vector $k_*(\mathbf{a}, \mathbf{s})$ whose
$i$-th entry is the cross-covariance between the $i$-th TD functional and
the latent $Q^*$ at the query point:
\begin{equation}
    [k_*(\mathbf{a}, \mathbf{s})]_i
    = k((\mathbf{a}_{i-1}, \mathbf{s}_{i-1}), (\mathbf{a}, \mathbf{s}))
    - \beta\, k((\mathbf{a}_i, \mathbf{s}_i), (\mathbf{a}, \mathbf{s}))
\end{equation}
The posterior mean and variance are then:
\begin{align}
    \hat{Q}^E_N(\mathbf{a}, \mathbf{s})
    &= \mu(\mathbf{a}, \mathbf{s})
    + \mathbf{k}_*(\mathbf{a}, \mathbf{s})^\top \, C^{-1} \,
      (\mathbf{Y} - \mathbf{M})
    \label{eq:dual_mean} \\
    \hat{\Sigma}^E_N((\mathbf{a}, \mathbf{s}), (\mathbf{a}, \mathbf{s}))
    &= \sigma^2_0
    - k_*(\mathbf{a}, \mathbf{s})^\top \, C^{-1} \,
      k_*(\mathbf{a}, \mathbf{s})
\end{align}
$C^{-1}(\mathbf{Y} - \mathbf{M})$ determines how much each past observation
revises the belief at a query point. $C^{-1}$ performs decorrelation
simultaneously across all observations. The inverse Gram matrix thus converts the
naive prior-based residuals $\mathbf{Y} - \mathbf{M}$ into the
sequentially conditioned innovations that the primal Kalman recursion
in \ref{def:experience_learning_update} would produce.
Dual form then allows us to reconstruct the posterior at any query point without ever
storing the posterior mean and variance functions themselves.

\paragraph{Worked example $t = 1$.}
The agent begins with prior $Q^* \sim \mathcal{GP}(\mu, k)$, takes action
$\mathbf{a}_0$ in state $\mathbf{s}_0$, receives utility
$\eta^E_1 = u(\mathbf{a}_0, \mathbf{s}_0)$, and transitions to state
$\mathbf{s}_1$. Define $\mathbf{x}^{\mathrm{dec}}_1 = (\mathbf{a}_0,
\mathbf{s}_0)$ and $\mathbf{x}^{\mathrm{out}}_1 = (\tilde{\mathbf{a}}_1,
\mathbf{s}_1)$, where $\tilde{\mathbf{a}}_1 = \arg\max_{\mathbf{a}}
\mu(\mathbf{a}, \mathbf{s}_1)$. The dual form stores these support points,
the scalar $\mathbf{Y} = (\eta^E_1)$, and the $1 \times 1$ inverse Gram
matrix $C^{-1} = (C_{11})^{-1}$, where
\begin{equation}
    C_{11} = k(\mathbf{x}^{\mathrm{dec}}_1, \mathbf{x}^{\mathrm{dec}}_1)
    - 2\beta\, k(\mathbf{x}^{\mathrm{dec}}_1, \mathbf{x}^{\mathrm{out}}_1)
    + \beta^2\, k(\mathbf{x}^{\mathrm{out}}_1, \mathbf{x}^{\mathrm{out}}_1)
    = \operatorname{Var}(\delta_1)
\end{equation}
is the prior variance of the first TD functional. The prior TD mean is
$M_1 = \mu(\mathbf{a}_0, \mathbf{s}_0) - \beta\,\mu(\tilde{\mathbf{a}}_1,
\mathbf{s}_1)$, so the residual $Y_1 - M_1 = \delta_1$. The weight is
$\alpha_1 = \delta_1 / C_{11}$. At any query point $(\mathbf{a}, \mathbf{s})$,
define $k_{*1} = k(\mathbf{x}^{\mathrm{dec}}_1, (\mathbf{a}, \mathbf{s}))
- \beta\, k(\mathbf{x}^{\mathrm{out}}_1, (\mathbf{a}, \mathbf{s}))$. Then:
\begin{equation}
    \hat{Q}^E_1(\mathbf{a}, \mathbf{s})
    = \mu(\mathbf{a}, \mathbf{s}) + \frac{k_{*1}}{C_{11}} \cdot \delta_1
    = \mu(\mathbf{a}, \mathbf{s}) + \alpha^E_1(\mathbf{a}, \mathbf{s}) \cdot \delta_1
\end{equation}
The ratio $k_{*1}/C_{11}$ is precisely the Kalman gain at step 1. 
The posterior variance is $\hat{\Sigma}^E_1 = \sigma^2_0 - k_{*1}^2 / C_{11}$.

\paragraph{Worked example $t = 2$.}
A second transition yields $\eta^E_2$, with new support points
$\mathbf{x}^{\mathrm{dec}}_2 = (\mathbf{a}_1, \mathbf{s}_1)$ and
$\mathbf{x}^{\mathrm{out}}_2 = (\tilde{\mathbf{a}}_2, \mathbf{s}_2)$.
The Gram matrix expands to $2 \times 2$. The new diagonal entry $C_{22}
= \operatorname{Var}(\delta_2)$ has the same structure as $C_{11}$.
The off-diagonal $C_{12} = \operatorname{Cov}(\delta_1, \delta_2)$ measures 
how correlated the two TD observations are (large when the observations are in nearby regions). 
Schur complement allows us to efficiently update the inverse Gram matrix without inverting the full $2 \times 2$ matrix.
\begin{equation}
    S = C_{22} - \frac{C_{12}^2}{C_{11}}
    = \operatorname{Var}(\delta_2 \mid \delta_1)
\end{equation}
The updated inverse is:
\begin{equation}
    C^{-1}_{\text{new}} = \begin{pmatrix}
        \frac{1}{C_{11}} + \frac{q^2}{S} & -\frac{q}{S} \\[4pt]
        -\frac{q}{S} & \frac{1}{S}
    \end{pmatrix}, \qquad q = \frac{C_{12}}{C_{11}}
\end{equation}
The weight vector $\boldsymbol{\alpha} = C^{-1}(\mathbf{Y} - \mathbf{M})$
now has two entries. The second is:
\begin{equation}
    \alpha_2
    = \frac{\delta^{\mathrm{prior}}_2 - q \cdot \delta^{\mathrm{prior}}_1}{S}
\end{equation}
where $\delta^{\mathrm{prior}}_i = \eta^E_i - M_i$ is the $i$-th TD error
evaluated against the original prior $\mu$. The numerator subtracts from
the raw residual of observation~2 the component linearly predicted by
observation~1 via the regression coefficient $q = C_{12}/C_{11}$. This is
the incremental surprise, or the Kalman innovation at step~2, which
the primal form obtains by measuring $\delta_2$ against the updated posterior
$\hat{Q}^E_1$ rather than $\mu$. The posterior at any query point is:
\begin{equation}
    \hat{Q}^E_2(\mathbf{a}, \mathbf{s})
    = \mu(\mathbf{a}, \mathbf{s})
    + k_{*1}\,\alpha_1 + k_{*2}\,\alpha_2
\end{equation}
The two forms yield identical posteriors: the primal form updates beliefs
sequentially so that each innovation is small; the dual form measures all
residuals against the prior and lets $C^{-1}$ decorrelate them. The pattern
generalizes to $N$ observations.